\scanpage{59}
espaces homogènes réguliers (ou ``distingués'').
\uline{Egalement} \unsure{?} :

\uline{Corollaire 3} Un polygône combinatoire de \unsure{rang $n$|rang|$n$-gone}
($2\leq n\leq\infty$) \deleted{muni} \uline{\unsure{de'un|de un|d'un} repère} est défini à
isom unique près -- et est \uline{can} \uline{isom}
à $\Pi_n$.

\uline{Corollaire 4}
\[
\begin{aligned}
\mathrm{Aut}(\Pi_n)&\simeq\unsure{D_n|\mathbb D_n}\quad(\text{can\textsuperscript{t}})\\
\mathrm{Aut}(\Pi)&\simeq D_n\quad(\text{non can\textsuperscript{t}})\ \text{si }\Pi\ n\text{-gône}
\end{aligned}
\]

\uline{NB} La restriction supplémentaire $\sigma_0x\ne\sigma_1x\quad\forall x\in E$
(en plus de $\sigma_0,\sigma_1$ ss pt fixe) sur les
\deleted{\unsure{espaces|orbites}} \inserted{$G_1$-ens} correspondants à des contours,
revient à l'\deleted{\unsure{élimination}} \inserted{exclusion} des orbites de la
forme $\deleted{\unsure{\Pi_2\simeq|\Pi_2=|\mathbb Z_2}}D_1$ -- \unsure{géom.|géos.} exclusion des
polygônes à 1 côté. On pourrait les
inclure, en n'exigeant pas $R\hookrightarrow S\times A$.

De ttes façons, = avec la restriction $n\geq 2$
pour les orbites, il semble qu'on soit encore
trop large dans la définition, si on compare
à celle des polyèdres \uline{combinatoires} -- où
on a exigé que \inserted{dans} l'ens ordonné associé
\deleted{deux} él. majorés aient un Sup. Ici \deleted{\unsure{ceti|cette|crois}}
l'exigence analogue reviendrait à exiger que
pour tte orbite on ait $n\geq3$ -- c'est à cette
% archive sheet number: 58
